\SecondArticle{9}
We shall complete the proof of the first theorem by showing that, when $Y$ and $Y'$ have no common divisor, $P=0$ certainly cannot hold. For this purpose, first suppose that two polynomials in the indeterminate $x$, namely $fx$ and $\varphi x$, have been found by the rules of Article 2, such that the identity
\[
fx\, .\,Y+\varphi x\, .\,Y'=1
\]
holds. We write it here as
\[
fx\, .\,\upsilon+\varphi x\, .\,\upsilon'=1+fx\, .\,(\upsilon-Y)+\varphi x\, .\,\frac{\mathrm d(\upsilon-Y)}{\mathrm dx}
\]
or, since we have
\[
\begin{aligned}
\upsilon'={}&(x-b)(x-c)(x-d)\, .\, .\, .\, .\\
&+(x-a)\, .\,\frac{\mathrm d[(x-b)(x-c)(x-d)\, .\, .\, .\, .]}{\mathrm dx}
\end{aligned}
\]
in the following form:
\[
\begin{aligned}
&\varphi x\, .\,(x-b)(x-c)(x-d)\, .\, .\, .\, .\\
&+\varphi x\, .\,(x-a)\, .\,\frac{\mathrm d[(x-b)(x-c)(x-d)\, .\, .\, .\, .]}{\mathrm dx}\\
&+fx\, .\,(x-a)(x-b)(x-c)(x-d)\, .\, .\, .\, .\\&=1+fx\, .\,(\upsilon-Y)+\varphi x\, .\,\frac{\mathrm d(\upsilon-Y)}{\mathrm dx}
\end{aligned}
\]
For brevity, we shall denote
\[
fx\, .\,(y-Y)+\varphi x\, .\,\frac{\mathrm d(y-Y)}{\mathrm dx}
\]
\SourceLeaf{050}
which is a polynomial in the indeterminates $x,l',l'',l'''$ etc.,
\[
\text{by}\quad F(x,l',l'',l'''\text{ etc.})
\]
Thus we have identically
\[
\begin{gathered}1+fx\, .\,(\upsilon-Y)+\varphi x\, .\,\frac{\mathrm d(\upsilon-Y)}{\mathrm dx}\\=1+F(x,\lambda',\lambda'',\lambda'''\text{ etc.})\end{gathered}
\]
We shall therefore have the identities [1]
\[
\begin{aligned}
&\varphi a\, .\,(a-b)(a-c)(a-d)\, .\, .\, .\, .\\&=1+F(a,\lambda',\lambda'',\lambda'''\text{ etc.})\\
&\varphi b\, .\,(b-a)(b-c)(b-d)\, .\, .\, .\, .\\&=1+F(b,\lambda',\lambda'',\lambda'''\text{ etc.})\\
&\varphi c\, .\,(c-a)(c-b)(c-d)\, .\, .\, .\, .\\&=1+F(c,\lambda',\lambda'',\lambda'''\text{ etc.})\\
&\text{etc.}
\end{aligned}
\]
Now suppose the product of all the expressions
\[
\begin{gathered}
1+F(a,l',l'',l'''\text{ etc.})\\
1+F(b,l',l'',l'''\text{ etc.})\\
1+F(c,l',l'',l'''\text{ etc.})\\
\text{etc.}
\end{gathered}
\]
which is a polynomial in the indeterminates $a,b,c$ etc., $l',l'',l'''$ etc., and is symmetric with respect to $a,b,c$ etc., to be represented by
\[
\psi(\lambda',\lambda'',\lambda'''\text{ etc.},\ l',l'',l'''\text{ etc.})
\]
Multiplying all the equations [1] then gives the new identity [2]
\[
\pi\varphi a\, .\,\varphi b\, .\,\varphi c\, .\, .\, .\, .=\psi(\lambda',\lambda'',\lambda'''\text{ etc.},\ \lambda',\lambda'',\lambda'''\text{ etc.})
\]
Further, since the product $\varphi a\, .\,\varphi b\, .\,\varphi c\, .\, .\, .\, .$ involves the indeterminates $a,b,c$ etc. symmetrically, we can clearly find a polynomial in the indeterminates $l',l'',l'''$ etc. which, under the substitutions $l'=\lambda',l''=\lambda'',l'''=\lambda'''$ etc., becomes $\varphi a\, .\,\varphi b\, .\,\varphi c\, .\, .\, .\, .$. Let $t$ be this polynomial. We then also have \mbox{identically [3]}
\[
pt=\psi(l',l'',l'''\text{ etc.},\ l',l'',l''')
\]
since these same substitutions $l'=\lambda',l''=\lambda'',l'''=\lambda'''$ etc. turn this equation into the identity [2].

Now the very definition of the function $F$ gives the identity
\SourceLeaf{051}
\[
F(x,L',L'',L'''\text{ etc.})=0
\]
Hence we also have identically
\[
\begin{gathered}
1+F(a,L',L'',L'''\text{ etc.})=1\\
1+F(b,L',L'',L'''\text{ etc.})=1\\
1+F(c,L',L'',L'''\text{ etc.})=1\\
\text{etc.}
\end{gathered}
\]
and therefore also identically
\[
\psi(\lambda',\lambda'',\lambda'''\text{ etc.},\ L',L'',L'''\text{ etc.})=1
\]
and thus identically [4]
\[
\psi(l',l'',l'''\text{ etc.},\ L'.\,L'',L'''\text{ etc.})=1
\]
Combining equations [3] and [4], and substituting $l'=L'$, $l''=L'',l'''=L'''$ etc., we shall therefore have [5]
\[
PT=1
\]
where $T$ denotes the value of the polynomial $t$ corresponding to those substitutions. Since this value must be a finite quantity, $P$ certainly cannot be $=0$. Q. E. D.

\SecondArticle{10}
It is now clear from the preceding results that any polynomial $Y$ in one indeterminate $x$ whose determinant is $=0$ can be decomposed into factors none of which has determinant 0. For, once the greatest common divisor of the polynomials $Y$ and $\frac{\mathrm dY}{\mathrm dx}$ has been found, the former is resolved into two factors. If one of these factors\HistoricalNote{11} again has determinant 0, it will be resolved into two factors in the same way. We shall continue in this manner until $Y$ is finally resolved into factors none of which has determinant 0.

It is also easy to see that among these factors into which $Y$ is resolved, at
\SourceLeaf{052}
least one must have the following property: among the factors of the number expressing its degree, the number two occurs no more often than among the factors of the number $m$ expressing the degree of the polynomial $Y$. That is, if we set $m=k\, .\,2^\mu$, with $k$ denoting an odd number, then among the factors of the polynomial $Y$ there is at least one of degree $k'\, .\,2^\nu$, where $k'$ is also odd, and either $\nu=\mu$ or $\nu<\mu$. This assertion follows at once from the fact that $m$ is the sum of the numbers expressing the degrees of the individual factors of $Y$.

\SecondArticle{11}
Before proceeding further, we shall explain an expression whose introduction is of great use in all investigations of symmetric functions, and which will also be very convenient for us. Suppose $M$ is a function of some of the indeterminates $a,b,c$ etc., and let $\mu$ be the number of them entering the expression $M$, without regard to any other indeterminates that $M$ may also involve. Permuting those $\mu$ indeterminates in every possible way, both among themselves and with the remaining $m-\mu$ of $a,b,c$ etc., produces from $M$ other expressions similar to $M$. In all there are
\[
m(m-1)(m-2)(m-3)\, .\, .\, .\, .\, .\,(m-\mu+1)
\]
expressions, including $M$ itself; we shall call their collection simply the \textit{collection of all $M$.} It is then clear what is meant by the sum of all $M$, the product of all $M$, etc. Thus, for example, $\pi$ is called the product of all $a-b$, $\upsilon$ the product of all $x-a$, and $\upsilon'$ the sum of all $\frac{\upsilon}{x-a}$, etc.

If $M$ is symmetric with respect to some of the $\mu$ indeterminates it contains, permutations of these among themselves do not change the function $M$. Thus every term occurs several times in the collection of all $M$: namely, $1\, .\,2\, .\,3\, .\, .\, .\, .\, .\,\nu$ times, if $\nu$ is the number of indeterminates with respect to which $M$ is symmetric. If $M$ is symmetric not only with respect to $\nu$ indeterminates, but also with respect to $\nu'$ others, and again with respect to $\nu''$ others, etc., then $M$ remains unchanged whether two of the first $\nu$ indeterminates are interchanged, or two of the second $\nu'$, or two of the third $\nu''$, etc. Thus
\[
1\, .\,2\, .\,3\, .\, .\, .\, .\, .\,\nu\, .\quad1\, .\,2\, .\,3\, .\, .\, .\, .\, .\,\nu'\, .\quad1\, .\,2\, .\,3\, .\, .\, .\, .\, .\,\nu''\text{ etc.}
\]
\SourceLeaf{053}
permutations always give identical terms. If we retain only one of these identical terms each time, we therefore obtain in all
\[
\frac{m(m-1)(m-2)(m-3)\, .\, .\, .\, .\, .\,(m-\mu+1)}{1\, .\,2\, .\,3\, .\, .\, .\, .\, .\,\nu\, .\,1\, .\,2\, .\,3\, .\, .\, .\, .\, .\,\nu'\, .\,1\, .\,2\, .\,3\, .\, .\, .\, .\, .\,\nu''\text{ etc.}}
\]
terms, whose collection we shall call the \textit{collection of all $M$ with repetitions excluded}, to distinguish it from the \textit{collection of all $M$ with repetitions included}. Unless an explicit qualification is given, we shall always understand repetitions to be included.

It is also easy to see that the sum of all $M$, or the product of all $M$, or more generally any symmetric function of all $M$, is always a symmetric function of the indeterminates $a,b,c$ etc., whether repetitions are included or excluded.

\SecondArticle{12}
With $u,x$ denoting indeterminates, consider now the product of all \mbox{$u-(a+b)x+ab$}, with repetitions excluded, and denote it by $\zeta$. Thus $\zeta$ is the product of the following \mbox{$\frac12m(m-1)$} factors:
\[
\begin{gathered}
u-(a+b)x+ab\\
u-(a+c)x+ac\\
u-(a+d)x+ad\\
\text{etc.}\\
u-(b+c)x+bc\\
u-(b+d)x+bd\\
\text{etc.}\\
u-(c+d)x+cd\\
\text{etc. etc.}
\end{gathered}
\]
Since this function involves the indeterminates $a,b,c$ etc. symmetrically, we can find a polynomial in the indeterminates $u,x,l',l'',l'''$ etc., to be denoted by $z$, which becomes $\zeta$ when $\lambda',\lambda'',\lambda'''$ etc. are substituted for the indeterminates $l',l'',l'''$ etc. Finally, denote by $Z$ the function of the indeterminates $u,x$ alone into which $z$ passes when the indeterminates $l',l'',l'''$ etc. are assigned the fixed values $L',L'',L'''$ etc.

These three functions $\zeta,z,Z$ can be regarded as polynomials of degree \mbox{$\frac12m(m-1)$} in the indeterminate $u$ with indeterminate coefficients. These
\SourceLeaf{054}
coefficients are
\[
\begin{aligned}
&\text{for }\zeta,\quad\text{functions of the indeterminates }x,a,b,c\text{ etc.}\\
&\text{for }z,\quad\text{functions of the indeterminates }x,l',l'',l'''\text{ etc.}\\
&\text{for }Z,\quad\text{functions only of the indeterminate }x.
\end{aligned}
\]
The individual coefficients of $z$ become those of $\zeta$ under the substitutions $l'=\lambda',l''=\lambda'',l'''=\lambda'''$ etc., and those of $Z$ under the substitutions $l'=L',l''=L'',l'''=L'''$ etc. What we have just said about the coefficients also holds for the determinants of the functions $\zeta,z,Z$. We shall now examine these more closely, with the aim of proving the following

\textsc{Theorem.} \textit{Whenever $P=0$ does not hold, the determinant of the function $Z$ certainly cannot be identically $=0$.}

